\begin{frame}[fragile]{}

\begin{itemize}
	\item As we learned the basics of information security, we can understand the applications that use the security tools.
\end{itemize}

\end{frame}

\begin{frame}[fragile]{Password}

\begin{columns}%
\begin{column}{0.7\textwidth}%
\begin{itemize}
	\item Servers do not store password in its original form.
	\item Instead, they create a hash and salt values from the password, and store them. 
	\item Therefore, even when the password files are hacked, hackers need to decrypt the original passwords from the hash values. 
\end{itemize}
\end{column}%
\begin{column}{0.3\textwidth}%
\begin{center}
\includegraphics[width=0.99\textwidth]{topics/SE_tools2/security/pic/crypto/hacked_password.png}
\end{center} 
\end{column}%
\end{columns}%
\end{frame}

\begin{frame}[fragile]{SSL}
\begin{itemize}
    \item SSL uses the hybrid 
	\item SSL certificates create a foundation of trust by establishing a secure connection.
	\item SSL aims to share a session key (symmetric/shared) securely between a browser and a web server. SSL uses a public key to share the session key.
	\item The shared session key is used for all communications between the browser and the server. 
\end{itemize}
\end{frame}

\begin{frame}[fragile]{How SSL Works?}

\begin{enumerate}
  \item Browser connects to a web server (website) secured with SSL (HTTPS). Browser requests that the server identify itself.
  \item Server sends a copy of its SSL Certificate, including the server's public key and a message.
\end{enumerate}

\begin{center}
\includegraphics[width=0.5\textwidth]{topics/SE_tools2/security/pic/ssl/ssl.png}
\end{center} 

\end{frame}

\begin{frame}[fragile]{}

\begin{enumerate}
\setcounter{enumi}{2}
  \item Browser checks the certificate root against a list of trusted CAs and that the certificate is unexpired, unrevoked, and that its common name is valid for the website that it is connecting to. 
  \item If the browser trusts the certificate, it creates, encrypts, and sends back a symmetric session key using the server's public key. Notice that the browser (client) and server uses hybrid encryption now. 
\end{enumerate}

\begin{center}
\includegraphics[width=0.5\textwidth]{topics/SE_tools2/security/pic/ssl/ssl.png}
\end{center} 

\end{frame}

\begin{frame}[fragile]{}

\begin{enumerate}
\setcounter{enumi}{4}
  \item Server decrypts the symmetric session key using its private key and sends back an acknowledgment encrypted with the session key to start the encrypted session.
  \item Server and Browser now encrypt all transmitted data with the session key.
\end{enumerate}

\begin{center}
\includegraphics[width=0.5\textwidth]{topics/SE_tools2/security/pic/ssl/ssl.png}
\end{center} 
\end{frame}

\begin{frame}[fragile]{TLS}

\begin{itemize}
	\item Transport Layer Security (TLS) is the successor protocol to SSL. TLS is an improved version of SSL. 
	\item It works in much the same way as the SSL, using encryption to protect the transfer of data and information. 
	\item The two terms are often used interchangeably in the industry, although SSL is still widely used.
\end{itemize}

\end{frame}

\begin{frame}[fragile]{HTTPS}

\begin{itemize}
	\item Hypertext Transfer Protocol Secure (HTTPS) is an extension of the Hypertext Transfer Protocol (HTTP). 
	\item It is used for secure communication over a computer network and is widely used on the Internet.
	\item In HTTPS, the communication protocol is encrypted using Transport Layer Security (TLS) or, formerly, Secure Sockets Layer (SSL). 
	\item The protocol is therefore also referred to as HTTP over TLS or HTTP over SSL.
\end{itemize}

\end{frame}

\begin{frame}[fragile]{SSH}
\begin{itemize}
	\item The Secure Shell Protocol (SSH) is a cryptographic network protocol for operating network services securely over an unsecured network.
	\item SSH operates on TCP port 22 by default (though this can be changed if needed). 
	\item The host (server) listens on port 22 (or any other SSH assigned port) for incoming connections. 
	\item It organizes the secure connection by authenticating the client and opening the correct shell environment if the verification is successful.
\end{itemize}
\end{frame}

\begin{frame}[fragile]{How Does SSH Work?}

\begin{itemize}
	\item When a client tries to connect to the server via TCP, the server presents the encryption protocols and respective versions that it supports. 
	\item If the client has a similar matching pair of a protocol and version, an agreement is reached, and the connection is started with the accepted protocol.
\end{itemize}
\end{frame}

\begin{frame}[fragile]{}

\begin{itemize}
	\item The server also uses an asymmetric public key which the client can use to verify the authenticity of the host.
	\item This is why the SSH client stores the public keys in \verb|~/.ssh/| for the verification.
	\item Once this is established, the two parties use what is known as a Diffie-Hellman Key Exchange Algorithm to create a symmetrical key.
\end{itemize}
\end{frame}

\begin{frame}[fragile]{SSH Authentication Using Password}

\begin{itemize}
	\item For authenticating users, most SSH users use a password.
	\item The user is asked to enter the username, followed by the password. 
	\item These credentials securely pass through the symmetrically encrypted tunnel, so there is no chance of them being captured by a third party.
\end{itemize}
\end{frame}

\begin{frame}[fragile]{Better Authentication Using Public Key}

\begin{itemize}
	\item We can generate and use SSH keys to enable simple and secure user authentication.
	\item We use `ssh-keygen -t rsa' to generate a pair of keys on the local machine. 
	\item The generated public key should be copied to the server using `ssh-copy-id user@serverip.'
	\item Then, we can log in to the remote server using ssh without giving the password. 
\end{itemize}

\end{frame}

\begin{frame}[fragile]{PGP}
\begin{itemize}
	\item Pretty Good Privacy (PGP) is an encryption program that provides cryptographic privacy and authentication for data communication.
	\item We have used JavaScript packages for information security.
	\item We also can use PGP tools for protecting digital information.
	\item We can use PGP as a library to be used for other programming languages - \href{https://www.openpgp.org/software/developer/}{https://www.openpgp.org/software/developer/}. 
\end{itemize}
\end{frame}

\begin{frame}[fragile]{GPG Installation}

\begin{itemize}
	\item GPG (GNU PGP) is the Open source implementation of PGP. 
	\item We can install PGP CLI programs using `brew' or `choco.'
	\item Use the name `gnupg' for its installation. When installed, the `gpg' is the command.
\end{itemize}

\begin{Command}{GPG}
\begin{lstlisting}[firstnumber=1, label=glabels, xleftmargin=0pt, basicstyle=\scriptsize]% 

brew install gnupg # Mac
choco install gnupg # PC
gnupg ... # need more arguments
\end{lstlisting}%   
\end{Command}

\end{frame}

\begin{frame}[fragile]{Create a public/private key pair}

\begin{itemize}
	\item `gpg --gen-key` is the command to make the key pair.
	\item Give your information with a paraphrase. I assume Sam (smcho) makes the key. 
	\item Then, a gpg directory is created to store the keys. 
\end{itemize}

\begin{Command}{}
\begin{lstlisting}[firstnumber=1, label=glabels, xleftmargin=0pt, basicstyle=\tiny]% 

gpg: /Users/smcho/.gnupg/trustdb.gpg: trustdb created
gpg: directory '/Users/smcho/.gnupg/openpgp-revocs.d' created
gpg: revocation certificate stored as '/Users/smcho/.gnupg/openpgp-revocs.d/F57AC82C7287D917A00C52AD951B59471C7CF6FE.rev'
public and secret key created and signed.

pub   ed25519 2022-03-15 [SC] [expires: 2024-03-14]
      F57AC82C7287D917A00C52AD951B59471C7CF6FE
uid                      smcho <chos5@nku.edu>
sub   cv25519 2022-03-15 [E] [expires: 2024-03-14]
\end{lstlisting}%   
\end{Command}
\end{frame}

\begin{frame}[fragile]{Make a Public Key}

\begin{itemize}
    \item Use \href{http://irtfweb.ifa.hawaii.edu/~lockhart/gpg/}{http://irtfweb.ifa.hawaii.edu/~lockhart/gpg/} as a reference. 
	\item Run `gpg --export -a "smcho" > public.key' to make a public key to share. 
	\item Students should use the correct user name when they used `pgp --gen-key' not "smcho". 
	\item Sam (smcho) can share the public.key to anyone. 
\end{itemize}

\end{frame}

\begin{frame}[fragile]{Importing a Shared Public Key}

\begin{itemize}
    \item Let's say Jim wants to send Sam (smcho) an encrypted file. 
    \item Jim receives Sam's public key (public.key). 
    \item Jim runs `gpg --import public.key' to import Sam's public key. 
    \item Jim can use `gpg --list-keys' to check the imported public keys. 
\end{itemize}

\end{frame}

\begin{frame}[fragile]{Encrypting a File}

\begin{itemize}
	\item Jim has a file (myfile.txt) with a secret text ("This is a secret message.") in it. 
	\item Run `gpg --recipient smcho --encrypt myfile.txt' to encrypt the file.
	\item Send the generated `myfile.txt.gpg' to Sam (smcho).
\end{itemize}

\end{frame}

\begin{frame}[fragile]{Decrypting a File}

\begin{itemize}
	\item Sam (smcho) can use `gpg -d myfile.txt.gpg' to recover the original file. 
	\item `gpg' requires to enter the `paraphrase' that Sam entered when he made a key pairs. 
\end{itemize}

\begin{Command}{}
\begin{lstlisting}[firstnumber=1, label=glabels, xleftmargin=0pt, basicstyle=\tiny]% 

smcho@mbp OneDrive-NorthernKentuckyUniversity> gpg -d myfile.txt.gpg
gpg: encrypted with cv25519 key, ID 46C764EADFBE878F, created 2022-03-15
      "smcho <chos5@nku.edu>"
This is a secret message.
\end{lstlisting}%   
\end{Command}

\end{frame}